\documentclass[12pt]{article}
\usepackage{seantex}

% --- Document Info ---
\title{Analysis II: Serie 09}
\author{Sean Findlay}
\date{\today}

% --- Start Document ---
\begin{document}

\maketitle

\begin{problem}{Volume}{}
    Calculate the volume of the region $B \subset \R^3$ enclosed by the surfaces $x^2 + y^2 + z^2 = 8$ and $2z = x^2 + y^2$.
\end{problem}

We would like to describe the region using cylindrical coordinates. Let $K$ be the parameter space of $\Phi$. Consider the transformation $\Phi: K \rightarrow B$:
$$
    \Phi(\rho, \vphi, z) = \begin{pmatrix} x \\ y \\ z \end{pmatrix} = \begin{pmatrix} \rho \cos \vphi \\ \rho \sin \vphi \\ z \end{pmatrix}
$$

We need:
$$
\begin{aligned}
    &(\rho \cos \vphi)^2 + (\rho \sin \vphi)^2 + z^2 = 8 \\
    &(\rho \cos \vphi)^2 + (\rho \sin \vphi)^2 = 2z \\
    &\implies 8 - z^2 = 2z \iff z^2 + 2z - 8 = 0 \iff (z + 4)(z - 2) = 0 \\
    &\implies z = 2 \text { or } z = -4
\end{aligned}
$$

Since $2z = x^2 + y^2$ and $x, y \in \R$, only the solution $z = 2$ is possible. So the two surfaces intersect at $z = 2$.

The two surfaces are rotationally symmetric around $z$, so we can have $0 \leq \vphi \leq 2\pi \iff \vphi \in [0, 2\pi]$.

For $\rho$ we have:
$$
\begin{aligned}
    (\rho \cos \vphi)^2 + (\rho \sin \vphi)^2 = 2 \cdot 2 = 4 &\iff \rho^2\left( \cos^2\vphi + \sin^2\vphi \right) = 4 \\
    &\iff \rho^2 \cdot 1 = 4 \iff \rho^2 = 4 \iff \rho = \pm 2
\end{aligned}
$$

We will only consider $\rho = 2$. So for the region enclosed by the surfaces, we need $\rho \in [0, 2]$

Finally, determine the full range of $z$ values. Given $\phi \in [0, 2\pi]$ and $\rho \in [0, 2]$, what values can $z$ take?
$$
\begin{aligned}
    x^2 + y^2 + z^2 = 8 &\iff z^2 = 8 - (x^2 + y^2) \iff z^2 = 8 - (\rho^2\cos^2\vphi + \rho^2\sin^2\vphi) \\
    &\iff z^2 = 8 - \rho^2(\cos^2\vphi + \sin^2\vphi) \iff z^2 = 8 - \rho^2 \\
    &\iff \pm \sqrt{z} = \pm \sqrt{8 - \rho^2}
\end{aligned}
$$

We will take the postive solution since we know $z = 2$ is the intersection. So at the boundary with this surface, $z = \sqrt{8 - \rho^2}$.

For the other surface, we have:
$$
\begin{aligned}
    2z = x^2 + y^2 &\iff 2z = \rho^2 \iff z = \frac{\rho^2}{2}
\end{aligned}
$$

So we have $z \in \left[\frac{\rho^2}{2}, \sqrt{8 - \rho^2}\right]$

Now, compute the Jacobian of $\Phi$ and its determinant:
$$
\begin{aligned}
    J \Phi(\rho, \vphi, z) = 
    \begin{pmatrix}
        \cos\vphi & -\rho\sin\vphi & 0 \\
        \sin\vphi & \rho\cos\vphi & 0 \\
        0 & 0 & 1 \\
    \end{pmatrix}
\end{aligned}
$$
$$
\det J\Phi(\rho,\vphi,z) = \det 
        \begin{pmatrix}
            \cos\vphi & -\rho\sin\vphi \\
        \sin\vphi & \rho\cos\vphi \\
        \end{pmatrix}
    = \rho\cos^2\vphi + \rho\sin^2\vphi = \rho
$$

So $|\det J \Phi| = \rho$. Since $J\Phi$ is a square matrix, we have that
$$
    \sqrt{\det\left(J\Phi^T J \Phi\right)} = |\det J \Phi|
$$



Given these bounds and the determinant, we can now compute the volume:

The volume is defined as:
$$
\begin{aligned}
    \mathrm{vol}_3(B) &= \int_B \sqrt{\det\left(J\Phi^T J \Phi\right)} \dd \lambda(\rho, \vphi, z) = \int_{0}^{2} \int_{0}^{2\pi} \int_{\rho^2/2}^{\sqrt{8 - \rho^2}} \rho \dd z \dd \vphi \dd \rho \\
    &= \int_{0}^{2\pi} \dd \vphi \int_{0}^{2} \int_{\rho^2/2}^{\sqrt{8 - \rho^2}} \rho \dd z \dd \rho = 2\pi \int_{0}^{2} \int_{\rho^2/2}^{\sqrt{8 - \rho^2}} \rho \dd z \dd \rho
\end{aligned}
$$

Compute the integral with respect to $z$:
$$
\begin{aligned}
    \int_{\rho^2/s}^{\sqrt{8 - \rho^2}} \rho \dd z &= \rho \Big[ z\Big]_{\rho^2/2}^{\sqrt{8 - \rho^2}} = \rho \left(\sqrt{8 - \rho^2} - \frac{\rho^2}{2}\right) = \rho\sqrt{8 - \rho^2} - \frac{\rho^3}{2}
\end{aligned}
$$

Substitute this into the original integral:
$$
\begin{aligned}
    \mathrm{vol}_3(B) &= 2\pi \int_{0}^{2} \int_{\rho^2/2}^{\sqrt{8 - \rho^2}} \rho \dd z \dd \rho = 2\pi \int_{0}^{2} \left(\rho\sqrt{8 - \rho^2} - \frac{\rho^3}{2} \right) \dd \rho \\
    &= 2\pi \left( \int_{0}^{2} \rho\sqrt{8 - \rho^2} \dd \rho - \int_{0}^{2} \frac{\rho^3}{2} \dd \rho \right) \\
\end{aligned}
$$

Compute the first term. Let $u := 8-\rho^2$. Then $\rho \dd \rho = -\frac{1}{2}\dd u$, so
$$
\begin{aligned}
    \int_{0}^{2} \rho\sqrt{8 - \rho^2} \dd \rho &= - \int_{4}^{8} \sqrt{u} \cdot -\frac{1}{2}\dd u = \frac{1}{2} \int_{4}^{8} u^\frac{1}{2} \dd u = \frac{1}{2} \left[ \frac{2}{3} u^\frac{3}{2} \right]_{4}^{8} \\
    &= \frac{1}{2} \cdot \frac{2}{3} \left( 8^\frac{3}{2} - 4^\frac{3}{2} \right) = \frac{1}{3}\left(\sqrt{512} - 8\right) = \frac{1}{3}\left(16 \sqrt{2} - 8\right)
\end{aligned}  
$$

Compute the second term:
$$
\begin{aligned}
    \int_{0}^{2} \frac{\rho^3}{2} \dd \rho &= \frac{1}{2} \int_{0}^{2} \rho^3 \dd \rho = \frac{1}{2} \left[ \frac{1}{4} \rho^4 \right]_0^2 = \frac{1}{8} \cdot 16 = 2
\end{aligned}
$$

So the volume is:
$$
\begin{aligned}
    &\mathrm{vol}_3(B) = 2\pi \left( \int_{0}^{2} \rho\sqrt{8 - \rho^2} \dd \rho - \int_{0}^{2} \frac{\rho^3}{2} \dd \rho \right) = 2\pi \left( \frac{1}{3}\left(16 \sqrt{2} - 8\right) - 2 \right)\\
    &\boxed{\mathrm{vol}_3(B) = \frac{4}{3}\pi\left(8 \sqrt{2} - 7\right)}
\end{aligned}
$$

\begin{problem}{Multiple integrals}{}
    \begin{enumerate}
        \item Let $D = [0, 2] \times [0, 1]$. Compute
        $$
            \iint_D \left( x^3 + 3x^2y +y^3 \right) \dd x \dd y
        $$

        \item Let $D \subset \R^2$ be the interior of the triangle with the vertices $(0, 0), (0, \pi), (\pi, \pi)$. Compute:
        $$
            \begin{aligned}
                \iint_D x\cos(x + y) \dd x \dd y
            \end{aligned}
        $$

        \item Let $D = \left\{ (x, y) \in \R^2 \ |\ x > 1,\ y > 1,\ x + y < 3 \right\}$. Compute:
        $$
            \iint_D \frac{1}{(x + y)^3} \dd x \dd y
        $$
    \end{enumerate}
\end{problem}

\begin{enumerate}
    \item Apply Fubini:
    $$
    \begin{aligned}
        \iint_{[0, 2] \times [0, 1]} \left( x^3 + 3x^2y +y^3 \right) \dd x \dd y &= \int_{0}^{2} \int_{0}^{1} \left( x^3 + 3x^2y +y^3 \right) \dd y \dd x
    \end{aligned}
    $$

    Compute the integral w.r.t. $y$:
    $$
    \begin{aligned}
        \int_{0}^{1} \left( x^3 + 3x^2y +y^3 \right) \dd y &= \left[x^3y + \frac{3}{2}x^2y^2 + \frac{1}{4}y^4\right]_{y = 0}^{y = 1} = x^3 + \frac{3}{2}x^2 + \frac{1}{4}
    \end{aligned}
    $$

    Substitute back into the original integral:
    $$
    \begin{aligned}
        \int_{0}^{2} \int_{0}^{1} \left( x^3 + 3x^2y +y^3 \right) \dd y \dd x &= \int_{0}^{2} x^3 + \frac{3}{2}x^2 + \frac{1}{4} \dd x = \left[ \frac{1}{4}x^4 + \frac{1}{2}x^3 + \frac{1}{4}x \right]_0^2 \\
        &= 4 + 4 + \frac{1}{2} = \boxed{\frac{17}{2}}
    \end{aligned}
    $$

    \item Compute the domain. The boundary consists of the line segments $y = \pi$, $x = 0$, $x = y$. So the domain is:
    $$
        D = \left\{ (x, y) \in \R^2 \ |\ 0 < x < y < \pi  \right\}
    $$
    
    Apply Fubini:
    $$
        \begin{aligned}
            \iint_D x\cos(x + y) \dd x \dd y &= \int_{0}^{\pi} \int_{x}^{\pi} x \cos(x + y) \dd y \dd x
        \end{aligned}
    $$

    Compute the integral w.r.t. $y$:
    $$
        \begin{aligned}
            \int_{x}^{\pi} x \cos(x + y) \dd y &= x \int_{x}^{\pi} \cos(x + y) \dd y = x \big[ \sin(x + y) \big]_{y = x}^{y = \pi} \\
            &= x\sin(x + \pi) - x\sin(2x) = -x\sin(x) -x\sin(2x)
        \end{aligned}
    $$

    Substitute back into the original integral:
    $$
        \begin{aligned}
            \int_{0}^{\pi} \int_{x}^{\pi} x \cos(x + y) \dd y \dd x &= \int_{0}^{\pi} \left(-x\sin(x) -x\sin(2x)\right) \dd x \\
            &= \int_{0}^{\pi} -x\sin(x) \dd x - \int_{0}^{\pi} x\sin(2x)\dd x \\
            &= \bigg[ x\cos(x) - \sin(x) \bigg]_0^\pi - \left[ -\frac{x}{2}\cos(2x) +\frac{1}{4} \sin(2x) \right]_0^\pi \\
            &= \pi\cos(\pi) - \sin(\pi) + \frac{\pi}{2}\cos(2\pi) - \frac{1}{4}\sin(2\pi) \\
            &= -\pi + \frac{\pi}{2} = \boxed{-\frac{\pi}{2}}
        \end{aligned}
    $$

    \item Apply Fubini:
    $$
    \begin{aligned}
        \iint_D \frac{1}{(x + y)^3} \dd x \dd y &= \int_1^2 \int_1^{3-x} \frac{1}{(x + y)^3} \dd y \dd x
    \end{aligned}
    $$

    Compute the integral w.r.t. $y$. Substitute: $u = x + y,\ \dd u = \dd y$:
    $$
    \begin{aligned}
        \int_1^{3-x} (x + y)^{-3} \dd y &= \int_{x + 1}^{3} u^{-3} \dd u = \left[-\frac{1}{2} u^{-2}\right]_{x+1}^3 \\
        &= -\frac{1}{2} \cdot \frac{1}{9} + \frac{1}{2}\cdot\frac{1}{(x+1)^2} = \frac{1}{2(x+1)^2} - \frac{1}{18}
    \end{aligned}
    $$

    Substitute back into the original integral:
    $$
    \begin{aligned}
        \int_1^2 \int_1^{3-x} \frac{1}{(x + y)^3} \dd y \dd x &= \int_1^2 \left( \frac{1}{2(x+1)^2} - \frac{1}{18} \right) \dd x \\
        &= \frac{1}{2} \int_1^2 (x+1)^{-2} \dd x - \frac{1}{18} \int_1^2 \dd x = 
    \end{aligned}
    $$

    Substitute: $u = x + 1$:
    $$
    \begin{aligned}
        \frac{1}{2} \int_1^2 (x+1)^{-2} \dd x - \frac{1}{18} \int_1^2 \dd x &= \frac{1}{2} \int_{2}^{3} u^{-2} \dd u - \frac{1}{18} = \frac{1}{2} \left[ -\frac{1}{u} \right]_{2}^{3} - \frac{1}{18} \\
        &= -\frac{1}{6} + \frac{1}{4} -\frac{1}{18} = \frac{-6 + 9 - 2}{36} \\
        &= \boxed{\frac{1}{36}}
    \end{aligned}
    $$
\end{enumerate}

\begin{problem}{Fubini's theorem}{}
    Compute the integrals:
    $$
        \int_0^{\sqrt{\pi}} \int_x^{\sqrt{\pi}} \sin(y^2) \dd y \dd x
    $$
    $$
        \int_{-1}^{1} \int_{|y|}^{1} (x + y)^2 \dd x \dd y
    $$
\end{problem}
\begin{enumerate}
    \item This integral is hard in the current order. If we could swap the order of the variables we integrate with respect to, we can make it easier. 
    
    The current domain is $D = \left\{(x, y) \in R^n \ |\ 0 < x < y < \sqrt{\pi}\right\}$


    We can apply Fubini's theorem:
    $$
    \begin{aligned}
        \int_0^{\sqrt{\pi}} \int_x^{\sqrt{\pi}} \sin(y^2) \dd y \dd x &= \int_0^{\sqrt{\pi}} \int_0^{y} \sin(y^2) \dd x \dd y
    \end{aligned}
    $$

    Now, the inner integral evaluates to:
    $$
    \begin{aligned}
        \int_0^{y} \sin(y^2) \dd x &= \sin(y^2) \int_0^{y} \dd x = \sin(y^2) \left[x\right]_0^y = y\sin(y^2)
    \end{aligned}
    $$

    Substitute back, let $u = y^2$, $\dd u = 2y \dd y \iff y \dd y = \frac{1}{2} \dd u$:
    $$
    \begin{aligned}
        \int_0^{\sqrt{\pi}} \int_0^{y} \sin(y^2) \dd x \dd y &= \int_0^{\sqrt{\pi}} y\sin(y^2) \dd y = \frac{1}{2} \int_0^{\pi} \sin(u) \dd u = \frac{1}{2}\Big[ -\cos(u) \Big]_0^\pi \\
        &= -\frac{1}{2} \cos(\pi) + \frac{1}{2} \cos(0) = \frac{1}{2} + \frac{1}{2} = \boxed{1}
    \end{aligned}
    $$

    \item We have the domain: $D = \left\{ (x, y) \in \R^2 \ |\ -1 < y < 1,\ |y| < x < 1 \right\}$ 
    
    The smallest that $|y|$ can be is $0$. So $(|y|, 1)$ is maximally $(0, 1)$. Using Fubini, we can write:
    $$
    \begin{aligned}
        \int_{-1}^{1} \int_{|y|}^{1} (x + y)^2 \dd x \dd y &= \int_{0}^{1} \int_{-x}^{x} (x + y)^2 \dd y \dd x
    \end{aligned}
    $$

    With the substitution $u = x + y$, $\dd u = \dd y$, the inner integral evaluates to:
    $$
        \begin{aligned}
            \int_{-x}^{x} (x + y)^2 \dd y &= \int_{0}^{2x} u^2 \dd u = \left[ \frac{1}{3} u^3 \right]_{0}^{2x} = \frac{8}{3}x^3
        \end{aligned}
    $$

    Substituting this back gives:
    $$
    \begin{aligned}
        \int_{0}^{1} \int_{-x}^{x} (x + y)^2 \dd y \dd x &= \int_{0}^{1} \frac{8}{3}x^3 \dd x = \frac{8}{3} \int_{0}^{1} x^3 \dd x = \frac{8}{3} \left[ \frac{1}{4} x^4 \right]_0^1 \\
        &= \frac{8}{3} \cdot \frac{1}{4} = \boxed{\frac{2}{3}}
    \end{aligned}
    $$
\end{enumerate}

\begin{problem}{Counterexample to Fubini}{}
    Let $f: [0, \infty)^2 \rightarrow \R$, defined by 
    $$
        f(x) =
        \begin{cases}
            e^{y-x} & x \geq y \geq 0, \\
            -e^{x-y} & 0 \leq x \leq y \\
        \end{cases}
    $$

    Compute the iterated integrals:
    $$
        \int_{0}^{\infty} \left( \int_{0}^{\infty} f(x, y) \dd x \right) \dd y,\quad \int_{0}^{\infty} \left( \int_{0}^{\infty} f(x, y) \dd y \right) \dd x
    $$

    and show that they have different values. Explain why this does not contradict Fubini's theorem.
\end{problem}

Compute the first integral:
$$
\begin{aligned}
    \int_{0}^{\infty} \left( \int_{0}^{\infty} f(x, y) \dd x \right) \dd y
\end{aligned}    
$$

We need to split the inner integral into its two pieces:
$$
\begin{aligned}
    \int_{0}^{\infty} f(x, y) \dd x &= \int_{0}^{y} -e^{x-y} \dd x + \int_{y}^{\infty} e^{y-x} \dd x
\end{aligned}  
$$

So we have:
$$
\begin{aligned}
    \int_{0}^{\infty} \left( \int_{0}^{\infty} f(x, y) \dd x \right) \dd y &= \int_{0}^{\infty} \left( \int_{0}^{y} -e^{x-y} \dd x + \int_{y}^{\infty} e^{y-x} \dd x \right) \dd y \\
\end{aligned}    
$$

The two parts evaluate to:
$$
\begin{aligned}
    \int_{0}^{y} -e^{x-y} \dd x &= \left[-e^{x-y}\right]_0^y = -e^{0} + e^{-y} = e^{-y} - 1 \\
    \int_{y}^{\infty} e^{y-x} \dd x &=  \left[ -e^{y-x} \right]_y^\infty = 0 - (-e^{0}) = 1
\end{aligned}
$$

So the integral is:
$$
\begin{aligned}
    \int_{0}^{\infty} \left( \int_{0}^{y} -e^{x-y} \dd x + \int_{y}^{\infty} e^{y-x} \dd x \right) \dd y &= \int_{0}^{\infty} \left( e^{-y} -1 + 1 \right) \dd y \\
    &= \int_{0}^{\infty} e^{-y} \dd y = \left[ -e^{-y} \right]_0^\infty \\
    &= 0 - (-1) = \boxed{1}
\end{aligned}   
$$

The second integral is:
$$
\begin{aligned}
    \int_{0}^{\infty} \left( \int_{0}^{\infty} f(x, y) \dd y \right) \dd x &= \int_{0}^{\infty} \left( \int_{0}^{x} e^{y-x} \dd y + \int_{x}^{\infty} -e^{x-y} \dd y \right) \dd x
\end{aligned} 
$$

The two parts evaluate to:
$$
\begin{aligned}
    \int_{0}^{x} e^{y-x} \dd y &= \left[ e^{y-x} \right]_0^x = e^0 - e^{-x} = 1 - e^{-x} \\
    \int_{x}^{\infty} -e^{x-y} \dd y &= \left[ e^{x-y} \right]_x^{\infty} = 0 - e^{0} = -1
\end{aligned}
$$

So the integral is:
$$
\begin{aligned}
    \int_{0}^{\infty} \left( \int_{0}^{x} e^{y-x} \dd y + \int_{x}^{\infty} -e^{x-y} \dd y \right) \dd x &= \int_{0}^{\infty} \left( 1 - e^{-x} - 1 \right) \dd x \\
    &= \int_{0}^{\infty} -e^{-x} \dd x = \left[ e^{-x} \right]_0^\infty\\
    &= 0 - e^0 = \boxed{-1}
\end{aligned} 
$$

So the two integrals are not equal. This is not a violation of Fubini's theorem, because of a requirement in the statement of the theorem (Theorem 13.47 from the lecture notes): The domain of integration must be bounded. In this problem, the domain is $[0, \infty)^2$, which cannot be contained within any finite bounding box. Therefore, the requirements for Fubini's Theorem are not met, and it cannot be applied.

\begin{problem}{Volume of the cone over a set}{}
    Let $\Omega \in \R^n$ be a bounded measurable set, $n \geq 1$. Consider the "cone over $\Omega$":
    $$
        C\Omega := \left\{ (x, t) \in \R^n \times \R \ |\ t \in [0, 1],\ x \in (1-t)\Omega \right\}
    $$

    Using Fubini's theorem and the homogeneity of $\mu_n$, show that
    $$
        \mu_{n+1}(C\Omega) = \frac{\mu_n(\Omega)}{n + 1}
    $$

    Use this result to compute the $n$-volume of the $n$-simplex:
    $$
        \mu_n(T_n) := \mu_n\left( \left\{ a_1e_1 + \dots + a_ne_n \ |\ a_i \in [0, 1],\ a_1 + \dots + a_n \leq 1 \right\} \right) = \frac{1}{n!}
    $$

    where $e_1,\dots,e_n$ denotes an orthonormal frame of $\R^n$.
\end{problem}

First, show that 
$$
\begin{aligned}
    \mu_{n+1}(C\Omega) = \frac{\mu_n(\Omega)}{n + 1}
\end{aligned}
$$

At any specific height $t$ of our cone, the "slice" at that height is precisely $(1-t)\Omega$. So, we can apply Fubini to integrate over these slices:
$$
\begin{aligned}
    \mu_{n+1}(C\Omega) = \int_{0}^{1} \mu_n((1-t)\Omega) \dd t
\end{aligned}
$$

By homogeneity of $\mu_n$, we have:
$$
\begin{aligned}
    \int_{0}^{1} \mu_n((1-t)\Omega) \dd t &= \int_{0}^{1} (1-t)^n \mu_n(\Omega) \dd t = \mu_n(\Omega) \int_{0}^{1} (1-t)^n \dd t \\
    &= \mu_n(\Omega) \int_{0}^{1} (1-t)^n \dd t = \mu_n(\Omega) \left[ \frac{-(1-t)^{n+1}}{n + 1} \right]_0^1 = \boxed{\frac{\mu_n(\Omega)}{n + 1}}
\end{aligned}
$$

Next, compute the volume
$$
    \mu_n(T_n) := \mu_n\left( \left\{ a_1e_1 + \dots + a_ne_n \ |\ a_i \in [0, 1],\ a_1 + \dots + a_n \leq 1 \right\} \right) = \frac{1}{n!}
$$

Let $a_n =: t$. Then,
$$
\begin{aligned}
    T_n &= \left\{ a_1e_1 + \dots + a_{n-1}e_{n-1} + te_n \ |\ a_i \in [0, 1],\ t \in [0, 1],\ a_1 + \dots + t \leq 1 \right\} \\
    &= \left\{ a_1e_1 + \dots + a_{n-1}e_{n-1} + te_n \ |\ a_i \in [0, 1],\ t \in [0, 1],\ a_1 + \dots + a_{n-1} \leq 1 - t \right\} \\
    &= \left\{ a_1e_1 + \dots + a_{n-1}e_{n-1} + te_n \ |\ a_i \in [0, 1],\ t \in [0, 1],\ \frac{a_1}{1 - t} + \dots + \frac{a_{n-1}}{1 - t} \leq 1 \right\} \\
\end{aligned}
$$

For all slices except $t \neq 1$ (which would be the single point at the top of the cone), we have
$$
    \frac{a_1}{1 - t} + \dots + \frac{a_{n-1}}{1 - t} \leq 1
$$

Because we have that $a_1 + \dots + a_{n-1} \leq 1 - t$, and all $a_i$ are positive, we know that $a_i \leq 1 -t$ for all $i = 1,\dots,n-1$. This implies $\frac{a_i}{1 - t} \leq 1$ for all $i = 1,\dots,n-1$. We also have that all $\frac{a_i}{1-t} \geq 0$. So the vector $(\frac{a_1}{1 - t}, \dots, \frac{a_{n-1}}{1 - t})$ is an element of $T_{n-1}$.

Let $x = a_1e_1 + \dots + a_{n-1}e_{n-1} = (a_1,\dots,a_{n-1})$. We have then that
$$
    (\frac{a_1}{1 - t}, \dots, \frac{a_{n-1}}{1 - t}) = \frac{1}{1-t}(a_1,\dots,a_{n-1}) = \frac{1}{1-t} x
$$

So we have 
$$
    \frac{1}{1-t} x \in T_{n-1} \iff x \in (1-t)T_{n-1}
$$

So we can rewrite:
$$
\begin{aligned}
    T_n &= \left\{ a_1e_1 + \dots + a_{n-1}e_{n-1} + te_n \ |\ a_i \in [0, 1],\ t \in [0, 1],\ a_1 + \dots + t \leq 1 \right\} \\
    &= \left\{ a_1e_1 + \dots + a_{n-1}e_{n-1} + te_n \ |\ a_i \in [0, 1],\ t \in [0, 1],\ \frac{a_1}{1 - t} + \dots + \frac{a_{n-1}}{1 - t} \leq 1 \right\} \\
    &= \left\{ x + te_n \ |\ t \in [0, 1],\ x \in (1-t)T_{n-1} \right\} = \left\{ (x, t) \in \R^{n-1} \times \R \ |\ t \in [0, 1],\ x \in (1-t)T_{n-1} \right\} \\
    &= CT_{n-1}
\end{aligned}
$$

So we can simply recursively apply the formula for the volume of a cone (as shown above):
$$
\begin{aligned}
    \mu_n(T_n) &= \mu_{n}\left(C(T_{n-1})\right) = \frac{1}{n} \mu_{n-1}(T_{n-1}) = \frac{1}{n}\cdot\frac{1}{n-1} \mu_{n-1}(T_{n-1}) = \dots = \frac{1}{n}\cdot \dots \cdot \frac{1}{2} \cdot \mu_1(T_1) \\
\end{aligned}
$$

Compute the value of $\mu_1(T_1)$:
$$
    \mu_1(T_1) = \mu_1({a_1e_1 \ |\ a_1 \in [0, 1],\ a_1 \leq 1 }) = \mu_1([0, 1]) = 1
$$

So we have:
$$
\begin{aligned}
    \mu_n(T_n) &= \frac{1}{n}\cdot \dots \cdot \frac{1}{2} \cdot \mu_1(T_1) \\
    &= \frac{1}{n \cdot \dots \cdot 2 \cdot 1} = \boxed{\frac{1}{n!}}
\end{aligned}
$$

\newpage

\begin{problem}{Gaussian integrals}{}
    Let $n \in \N$ and $A \in \mathrm{M}_{n\times n}(\R)$ be a symmetric positive definite matrix. Show that
    $$
        \int_{\R^n} e^{-\langle Ax, x \rangle}\dd x = \frac{\pi^{n/2}}{\sqrt{\det(A)}}
    $$
\end{problem}

Prove first the specific case where $A$ is a diagonal matrix. Let
$$
    A =: \begin{pmatrix}
        \lambda_1 & \dots & 0 \\
        \vdots & \ddots & \vdots \\
        0 & \dots & \lambda_n \\
    \end{pmatrix}
$$

Then for $x \in \R^n,\ x =: \begin{pmatrix} x_1 \\ \vdots \\ x_n \end{pmatrix}$, we have
$$
    \langle Ax, x \rangle = (Ax)^T x = \begin{pmatrix} \lambda_1x_1\ \dots\ \lambda_nx_n \end{pmatrix} \begin{pmatrix} x_1 \\ \vdots \\ x_n \end{pmatrix} = \lambda_1x_1^2 + \dots +\lambda_nx_n^2 
$$

Then, the integral is:
$$
\begin{aligned}
    \int_{\R^n} e^{-\langle Ax, x \rangle}\dd x &= \int_{\R^n} e^{-(\lambda_1x_1^2 + \dots +\lambda_nx_n^2) }\dd x = \int_{\R^n} e^{-\lambda_1x_1^2} \cdot\ \dots\ \cdot e^{-\lambda_nx_n^2 }\dd x \\
\end{aligned}
$$

Apply Fubini:
$$
\begin{aligned}
    \int_{\R^n} e^{-\lambda_1x_1^2} \cdot\ \dots\ \cdot e^{-\lambda_nx_n^2 }\dd x &= \int_\R \dots \int_\R e^{-\lambda_1x_1^2} \cdot\ \dots\ \cdot e^{-\lambda_nx_n^2 } \dd x_n \dots \dd x_1 \\
    &= \int_\R e^{-\lambda_1x_1^2} \dd x_1 \cdot \dots \cdot \int_\R e^{-\lambda_nx_n^2} \dd x_n \\
\end{aligned}
$$


We can use the following integral to compute this:
$$
\begin{aligned}
    \int_\R e^{-x^2} \dd x = \sqrt{\pi} \iff \int_\R e^{-\lambda x^2} \dd x &= \int_\R e^{-(\sqrt{\lambda}x)^2} \dd x = \int_\R e^{-u^2} \frac{1}{\sqrt{\lambda}} \dd u \\
    &= \frac{1}{\sqrt{\lambda}} \int_\R e^{-u^2} \dd u = \sqrt{\frac{\pi}{\lambda}}
\end{aligned}
$$

Apply this:
$$
\begin{aligned}
    \int_\R e^{-\lambda_1}e^{x_1^2} \dd x_1 \cdot \dots \cdot \int_\R e^{-\lambda_n}e^{x_n^2} \dd x_n &= \sqrt{\frac{\pi}{\lambda_1}} \dots \sqrt{\frac{\pi}{\lambda_n}} = \sqrt{\frac{\pi^n}{\lambda_1 \dots \lambda_n}} \\
    &= \frac{\pi^{n/2}}{\sqrt{\det(A)}}
\end{aligned}
$$

So the statement holds for diagonal matrices. Now, show the general case where $A$ is only symmetric and positive definite. 

Let $A$ be a symmetric positive definite matrix. By the spectral theorem, $A$ is diagonalisable. So there exists an orthogonal $Q \in \mathrm{GL}_n(\R)$ such that $D := Q^{-1}AQ$ is a diagonal matrix. Because $D$ is diagonal, it holds that
$$
    \int_{\R^n} e^{-\langle Dx, x \rangle}\dd x = \int_{\R^n} e^{-\langle (Q^{-1}AQ)x, x \rangle}\dd x = \frac{\pi^{n/2}}{\sqrt{\det(D)}}
$$

We want to show that:
$$
    \int_{\R^n} e^{-\langle (Q^{-1}AQ)x, x \rangle}\dd x = \int_{\R^n} e^{-\langle Ax, x \rangle}\dd x
$$

Given that $Q$ is orthogonal, we have $Q^{-1} = Q^T$, so
$$
\begin{aligned}
    \langle (Q^{-1}AQ)x, x \rangle &= \langle (Q^TAQ)x, x \rangle = ((Q^TAQ)x)^T x = x^T(Q^TAQ)^T x = x^T Q^T A Q x = (Qx)^T A (Qx) \\
    &= (Qx)^T A^T (Qx) = (A(Qx))^T (QX) = \langle A(Qx), Qx \rangle
\end{aligned}
$$

So we can use the substitution $w := Qx$, $\dd w = |\det(Q)| \dd x = |\pm 1| \dd x = \dd x$ and rewrite:
$$
    \int_{\R^n} e^{-\langle (Q^{-1}AQ)x, x \rangle}\dd x = \int_{\R^n} e^{-\langle A(Qx), Qx \rangle}\dd x = \int_{\R^n} e^{-\langle Aw, w \rangle}\dd w
$$

Because $w$ is just now just a symbol, we have shown that
$$
    \int_{\R^n} e^{-\langle (Q^{-1}AQ)x, x \rangle}\dd x = \int_{\R^n} e^{-\langle Ax, x \rangle}\dd x
$$

We also have that
$$
\begin{aligned}
    \det(D) &= \det(Q^{-1}AQ) = \det(Q^{-1})\det(A)\det(Q) = \det(Q^T)\det(A)\det(Q) = \det(Q)^2\det(A) \\
    &= \det(A)
\end{aligned}
$$

So
$$
    \frac{\pi^{n/2}}{\sqrt{\det(D)}} = \frac{\pi^{n/2}}{\sqrt{\det(A)}}
$$

Because $\int_{\R^n} e^{-\langle (Q^{-1}AQ)x, x \rangle}\dd x = \int_{\R^n} e^{-\langle Ax, x \rangle}\dd x$ and $\frac{\pi^{n/2}}{\sqrt{\det(D)}} = \frac{\pi^{n/2}}{\sqrt{\det(A)}}$, we have shown that the statement also holds for symmetric positive definite matrices. So the statement holds. \qed

\begin{problem}{Layer-cake formula}{}
    Let $f: \R^n \rightarrow [0, \infty)$ be a continuous function which vanishes identically outside a compact set, and let $p \geq 1$. Using Fubini's theorem, show the Layer-cake formula:
    $$
        \int_\R f(x)^p \dd x = p \int_{0}^{\infty} t^{p-1} \mu_n \left(\left\{ x \in \R^n \ |\ f(x) > t \right\}\right) \dd t
    $$

    Find a similar formula for the integral of
    $$
        \int_\R \Phi(f(x)) \dd x = \int_{0}^{\infty} \dots \mu_n\left( \left\{x \in \R^n \ |\ f(x) > t\right\}\right) \dd t
    $$

    where $\Phi \in C^1(\R)$ is any function such that $\Phi(0) = 0$.
\end{problem}

$f(x)$ is just some non-negative real number. So, by the fundamental theorem of calculus,
$$
    f(x)^p = \int_{0}^{f(x)} p \cdot t^{p-1} \dd t
$$

Rewrite:
$$
\begin{aligned}
    \int_\R f(x)^p \dd x &= \int_{-\infty}^{\infty} \int_{0}^{f(x)} p \cdot t^{p-1} \dd t \dd x
\end{aligned}
$$

Define $A := [0, f(x)]$. Then, let $\mathds{1}_A$ be the indicator function for the set $A$:
$$
    \mathds{1}_A(t) = 
    \begin{cases}
        1 & t \in A, \\
        0 & t \notin A
    \end{cases}
$$

Then, we can extend the bounds of the inner integral:
$$
\begin{aligned}
    \int_\R f(x)^p \dd x &= \int_{-\infty}^{\infty} \int_{0}^{f(x)} p \cdot t^{p-1} \dd t \dd x = \int_{-\infty}^{\infty} \int_{0}^{\infty} \mathds{1}_A(t) \cdot p \cdot t^{p-1} \dd t \dd x
\end{aligned}
$$

Now, we can apply Fubini:
$$
\begin{aligned}
    \int_{-\infty}^{\infty} \int_{0}^{\infty} \mathds{1}_A(t) \cdot p \cdot t^{p-1} \dd t \dd x &= \int_{0}^{\infty} \int_{-\infty}^{\infty} \mathds{1}_A(t) \cdot p \cdot t^{p-1} \dd x \dd t \\
    &= \int_{0}^{\infty} p t^{p-1} \left( \int_{-\infty}^{\infty} \mathds{1}_A(t) \dd x \right) \dd t \\
\end{aligned}
$$

Notice that the inner integral is $1$ whenever $f(x) > t$, because then we have $t \in [0, f(x)]$. So we are basically computing the measure of the following set:
$$
    \int_{-\infty}^{\infty} \mathds{1}_A(t) \dd x = \mu_n\left(\left\{x \in \R^n \ |\ f(x) > t\right\}\right)
$$

Substitute this back in order to get
$$
\begin{aligned}
    \int_\R f(x)^p \dd x &= \int_{0}^{\infty} p t^{p-1} \left( \int_{-\infty}^{\infty} \mathds{1}_A(t) \dd x \right) \dd t = p \int_{0}^{\infty} t^{p-1} \mu_n\left(\left\{x \in \R^n \ |\ f(x) > t\right\}\right) \dd t
\end{aligned}
$$

which is what we wanted to show.

Next, find a formula for the integral
$$
    \int_\R \Phi(f(x)) \dd x
$$

where $\Phi \in C^1(\R)$ with $\Phi(0) = 0$. This means $\Phi'$ exists and is continuous. So, we can apply the fundamental theorem of calculus again:
$$
    \Phi(z) = \int_0^{z} \Phi'(t) \dd t
$$

Plug in $f(x)$ (which is just a non-negative real number):
$$
    \Phi(f(x)) = \int_0^{f(x)} \Phi'(t) \dd t
$$

Substitute into our integral:
$$
    \int_\R \Phi(f(x)) \dd x = \int_{-\infty}^{\infty} \int_0^{f(x)} \Phi'(t) \dd t \dd x
$$

Let $A := [0, f(x)]$ and let $\mathds{1}_A$ be the indicator function for this set. Then we can rewrite and apply Fubini like before:
$$
\begin{aligned}
    \int_{-\infty}^{\infty} \int_0^{f(x)} \Phi'(t) \dd t \dd x &= \int_{-\infty}^{\infty} \int_{0}^{\infty} \mathds{1}_A(t) \cdot \Phi'(t) \dd t \dd x = \int_{0}^{\infty} \int_{-\infty}^{\infty} \mathds{1}_A(t) \cdot \Phi'(t) \dd x \dd t \\
    &= \int_{0}^{\infty} \Phi'(t) \left( \int_{-\infty}^{\infty} \mathds{1}_A(t) \dd x \right) \dd t \\
    &= \int_{0}^{\infty} \Phi'(t) \cdot \mu_n\left(\left\{ x \in \R^n \ |\ f(x) > t \right\}\right) \dd t
\end{aligned}
$$

which is a formula for the integral we wanted. \qed

\end{document}